\chapter{Related Work} \label{Related Work}

\section{Image Editing Models and the State of the Art}

\label{Image Editing Models}

The landscape of instruction-guided image editing has evolved rapidly, with models spanning multiple architectural paradigms now competing at the frontier. InstructPix2Pix \cite{brooks2023instructpix2pix} established the foundational paradigm of paired instruction-based editing, training a diffusion model on synthetically generated editing pairs to enable text-guided image modification. While this work did not address text-in-image editing, it demonstrated that instruction-conditioned diffusion models could learn meaningful editing operations from paired data, a principle that underpins much subsequent work.

On the diffusion side, Qwen-Image \cite{wu2025qwenimage} represents the current open-source state of the art for image generation and editing, with particular strength in text rendering. Qwen-Image introduces a dual-encoding mechanism that feeds the source image through both a vision-language model and a VAE encoder, enabling the editing module to balance semantic consistency with visual fidelity. Its progressive training strategy for text rendering, scaling from simple to complex textual inputs, has yielded strong performance on both alphabetic and logographic scripts. For scene text editing specifically, TextCtrl \cite{zeng2024textctrl} addresses content replacement through disentangled style and glyph structure priors, achieving high fidelity in changing what text says. However, neither Qwen-Image nor TextCtrl targets quantitative spatial editing of text, such as repositioning or resizing text elements by precise amounts.

On the autoregressive side, VAREdit \cite{mao2026varedit} reframes instruction-guided image editing as a next-scale prediction problem within the Visual Autoregressive (VAR) framework \cite{tian2024var}. VAREdit introduces a Scale-Aligned Reference module that injects scale-matched conditioning from the source image, and demonstrates that the compositional, token-by-token generation of autoregressive models naturally avoids the entanglement problems that plague diffusion-based approaches. On standard editing benchmarks, VAREdit outperforms leading diffusion methods in instruction adherence while achieving faster inference. This architectural distinction between diffusion and autoregressive approaches becomes central to our comparative analysis in this work.

Among commercial systems, Gemini 3 Pro Image (marketed as Nano Banana Pro) \cite{google2025nanobanana} has emerged as arguably the strongest practical model for image editing, including text manipulation tasks. According to its model card, it is based on the Gemini 3 Pro architecture, a sparse mixture-of-experts (MoE) transformer with native multimodal capabilities, placing it in the autoregressive family rather than the diffusion paradigm \cite{google2025promodelcard}. Despite leading on internal Elo benchmarks across text rendering and editing capabilities, even Nano Banana Pro's own documentation acknowledges persistent difficulties with spatial localisation and small text rendering \cite{google2025proimagemodelcard}. These self-reported limitations align with the broader challenges in quantitative text editing that motivate our work.


\section{Architectural Limitations for Quantitative Editing}

\label{Architectural Limitations}

The difficulty of quantitative text editing is not merely a matter of insufficient training data; it is rooted in fundamental architectural constraints of current generative models. Understanding these constraints is essential for interpreting why certain failure modes persist and for motivating architectural comparisons.

Diffusion-based models generate images through an iterative denoising process that operates globally over the entire image at each step. This global denoising inherently entangles the edited region with the full image context, making it difficult to perform spatially precise, localized modifications without introducing unintended changes elsewhere \cite{mao2026varedit}. For text regions specifically, Shi et al. \cite{shi2025wordcon} have shown that attention map misalignment is more prominent in text rendering than in common object generation, with the decoupling between words used for text rendering being less thorough than for common objects. This indicates that current diffusion methods lack effective mechanisms for precise text positioning, making word-level controllability particularly challenging. These properties suggest that diffusion architectures may systematically struggle with tasks requiring pixel-level spatial precision, such as repositioning text to exact coordinates or scaling text elements by a specified factor while preserving surrounding content.

Visual Autoregressive models face a different but equally consequential set of challenges. The VAR architecture \cite{tian2024var} generates images through next-scale prediction, progressively building from coarse, low-resolution representations to fine, high-resolution detail. While this coarse-to-fine approach produces high-quality images and avoids the entanglement problem of diffusion, it introduces complications for editing tasks. Mao et al.\ \cite{mao2026varedit} identify a scale-mismatch problem in which source image features at the finest scale cannot effectively guide prediction at coarser scales, and propose a Scale-Aligned Reference module to inject scale-matched conditioning as a remedy. That the multi-scale conditioning required dedicated architectural intervention suggests that the interaction between scales is a critical challenge for VAR-based editing.

More broadly, recent work has identified fundamental limitations in the autoregressive image generation pipeline that compound the difficulty of pixel-precise tasks. He et al.\ \cite{he2025rear} demonstrate that a core bottleneck lies in generator-tokenizer inconsistency: the token sequences produced by the autoregressive generator may fall outside the distribution that the tokenizer's decoder can faithfully reconstruct, causing token-level errors to propagate into perceptual and pixel-level degradation in the decoded image. Their analysis shows that even when a generated token sequence appears reasonable by language-modelling metrics such as perplexity, the decoded image can still suffer significant fidelity loss, indicating that the error is not merely one of prediction quality but of a structural mismatch between generation and decoding. This inconsistency is particularly concerning for quantitative editing, where even small deviations at the token level can translate into measurable spatial or colour errors in the output. Empirical evidence from the commercial domain reinforces these concerns: the Gemini 3 Pro Image model card \cite{google2025proimagemodelcard} acknowledges that even this leading autoregressive system exhibits occasional confusion around spatial localisation and poor rendering of small text, while Zuo et al.\ \cite{zuo2025nanobananapro} document a broader pattern in which generative models achieve strong subjective visual quality yet lag behind specialist models on reference-based quantitative metrics, attributing this gap to the inherent stochasticity of generative models and their difficulty maintaining strict pixel-level consistency.

Together, these findings suggest that the pixel-fidelity challenges of autoregressive models stem from multiple interacting sources, scale mismatches in the coarse-to-fine pipeline, generator-tokenizer inconsistency and error accumulation during sequential decoding, and the stochastic nature of generative reconstruction, whose combined effect on precise quantitative transformations remains poorly understood. Characterizing how these limitations manifest as distinct failure modes on quantitative text editing is a goal of our work.

\section{Evaluation of Text Editing in Images}

\label{Evaluation}

Evaluating quantitative text editing requires benchmarks and metrics capable of measuring precise, measurable changes, a requirement that existing evaluation frameworks only partially fulfill.

Raji et al. \cite{raji2021everything} argue that widely adopted AI benchmarks suffer from fundamental construct validity problems: they claim to measure general capabilities but are in practice finite, context-specific artifacts whose composite scores can obscure more than they reveal. Most relevant to our work is their observation that aggregate metrics can mask failures on specific sub-tasks, and their recommendation that evaluation should move toward diagnostic approaches, including test suites, disaggregated analysis that isolates particular failure modes across subgroups, and comparative methods such as ablation testing that reveal how different system designs respond to the same challenges.

TextEditBench \cite{gui2025texteditbench} is the most directly relevant benchmark to our work, as it explicitly evaluates text editing within images across multiple dimensions including content modification and spatial transformation. Critically, TextEditBench documents that current state-of-the-art models still struggle significantly with quantitative aspects such as text relocalization and spatial adjustment. However, viewed through the lens of Raji et al.'s critique, TextEditBench exhibits some of the same limitations common to broad benchmarks. Its evaluation includes reasoning-intensive scenarios and composite metrics that can conflate multiple failure modes, making it difficult to isolate whether a model fails because it cannot execute a spatial transformation, because it does not act without explicit instruction, or because it cannot interpret a complex instruction. In the terminology of Raji et al., the benchmark's aggregate reporting risks obscuring architecture-specific weaknesses behind a single performance number.

Our evaluation framework addresses this directly by following Raji et al.'s recommendation toward diagnostic, disaggregated evaluation. Rather than measuring broad editing competence through composite scores, we test atomic, isolated operations: did the text move, and how close to the ground truth position? Did the size change, and by how much does the result deviate from the target? This design allows us to pinpoint specific failure modes, including spatial relocation accuracy, scaling precision, color fidelity, and unintended modifications, rather than conflating them, yielding the kind of fine-grained, interpretable signal that Raji et al. argue general benchmarks structurally cannot provide.

This need for decomposed evaluation extends beyond spatial precision to style preservation. STELLAR \cite{seo2025stellar} identifies a similar evaluation gap in scene text editing: standard metrics fail to capture whether visual style has been faithfully preserved after an edit. Their proposed Text Appearance Similarity (TAS) metric uses a trained style encoder to independently measure font, color, and background similarity between source and edited images, enabling reference-free style evaluation. This is directly relevant to quantitative text editing, where a model that correctly resizes text but inadvertently alters its font has failed in a way that composite metrics would obscure.

The metric misalignment identified by Zuo et al. \cite{zuo2025nanobananapro}, where generative models score poorly on quantitative metrics despite producing perceptually strong outputs, also warrants consideration in our evaluation design. Their argument that ``ground truth is not the unique optimum for generative repair'' holds for many low-level vision tasks where multiple plausible outputs exist. However, quantitative text editing represents a setting where this concern is substantially reduced: if the instruction is ``move the text 50 pixels to the right,'' there is a single correct spatial outcome, and pixel-level metrics against a ground-truth rendering provide a meaningful and fair assessment.

More broadly, several recent benchmarks address image editing without specifically targeting text: EDITVAL \cite{basu2023editval} and HYPE-EDIT-1 \cite{chan2026hypeedit1} evaluate general image editing quality, with HYPE-EDIT-1 notably identifying typography and layout as underserved areas in its future work. However, none of these benchmarks provide the paired ground-truth data or fine-grained spatial metrics needed to evaluate quantitative text editing operations.

